% SGA 3，第 XI 篇；根据英文重建源翻译。
% 范围：第 5 节全部，从标题至注 5.11；止于第 6 节标题之前。
% 权威来源：Polo--Gille Expo11.pdf，组件局部页 25--29；合卷页 747--751；印刷／原文页边页 165--171。
% Jacob Reinhold 的 CC BY 4.0 Markdown 谱系仅用于措辞比较；文字、公式、编号、注释及图均以 Polo--Gille PDF 为准。
% 精确续接／终点：第 6 节标题，Expo11 局部页 29／合卷页 751／印刷页 171。

\SGARefTarget{sga3:xi:section:5}\section*{5. 可表性定理的首批推论}\addcontentsline{toc}{section}{5. 可表性定理的首批推论}

设 \(H\)、\(G\) 如 \SGARefLink{sga3:xi:corollary:4:2}{4.2} 中所述。置 \SGARefTarget{sga3:xi:display:5:hom-scheme} \[
  M=\operatorname{Hom}_{S\text{-}\mathrm{gr}}(H,G),
\]，由 \SGARefLink{sga3:xi:corollary:4:2}{4.2}，它是 \(S\) 上分离的平滑 \(S\)-预概形。注意 \(G\) 通过 \SGARefTarget{sga3:xi:display:5:conjugation-action} \[
  (g,u)\longmapsto\operatorname{int}(g)\circ u.
\] 作用于函子 \(M=\operatorname{Hom}_{S\text{-}\mathrm{gr}}(H,G)\)。因此存在典范态射 \SGARefTarget{sga3:xi:display:5:x} \[
  \tag{x}
  G\times_S M\longrightarrow M\times_S M,
\]，其分量为 \SGARefLink{sga3:xi:display:5:conjugation-action}{上述态射} \(G\times_S M\to M\) 及第二投影 \(G\times_S M\to M\)。

\medskip\noindent\SGARefTarget{sga3:xi:corollary:5:1}\textbf{推论 5.1。} --- \textit{上述态射 \(\mathrm{\SGARefLink{sga3:xi:display:5:x}{(x)}}\) 平滑。}

这由 \SGARefLink{sga3:xi:corollary:4:2}{4.2} 和 \SGARefLink{sga3:xi:corollary:2:3}{2.3} 得出。此陈述等价于下述陈述。

\medskip\noindent\SGARefTarget{sga3:xi:corollary:5:2}\textbf{推论 5.2。} --- \textit{设 \(u_1,u_2:H\rightrightarrows G\) 为两个 \(S\)-群同态。则 \(G\) 的子函子 \(\operatorname{Transp}(u_1,u_2)\)（参见 \SGARefLink{sga3:xi:corollary:2:4}{2.4}）可由 \(G\) 的闭子预概形表示，且在 \(S\) 上平滑。}

只须证明 \(\operatorname{Transp}(u_1,u_2)\to G\) 确为闭浸入。这是因为它是一对态射 \(G\rightrightarrows M\) 的等化子，具体地，这对态射为 \[
  g\longmapsto\operatorname{int}(g)u_1
  \qquad\text{且}\qquad
  g\longmapsto u_2,
\]，其中后者为“常”态射，并且 \(M\) 在 \(S\) 上分离。特别地：

\medskip\noindent\SGARefTarget{sga3:xi:corollary:5:3}\textbf{推论 5.3。} --- \textit{设 \(u:H\to G\) 为 \(S\)-群态射。则 \[
  \operatorname{Centr}_G(u)=\operatorname{Transp}(u,u)
\] 可由 \(G\) 的闭群子预概形表示，且在 \(S\) 上平滑。此外，\(G\) 对 \(\operatorname{Centr}_G(u)\) 的商可由 \(M\) 的开子预概形表示。}

还须证明最后一个断言。由 \SGARefLink{sga3:xi:corollary:5:2}{5.2}，从 \(G\) 到 \(M\) 的态射 \SGARefTarget{sga3:xi:display:5:orbit-map} \[
  g\longmapsto\operatorname{int}(g)\circ u
\] 平滑且有限型，因而为开态射（EGA IV，6.6）。以 \(U\) 表示其像，赋予从 \(M\) 诱导的结构。诱导态射 \(G\to U\) 平滑、满射且有限型，因此为忠实平坦拟紧拓扑的覆盖。此外，显然 \SGARefLink{sga3:xi:display:5:orbit-map}{上述态射} \(G\to M\) 使 \(G\) 成为 \(\operatorname{Centr}_G(u)_M\) 作用下的形式主齐性层。因此层 \(G/\operatorname{Centr}_G(u)\) 由 \(U\) 表示。

\medskip\noindent\SGARefTarget{sga3:xi:corollary:5:4}\textbf{推论 5.4。} --- \textit{设 \(u_1,u_2:H\rightrightarrows G\) 为两个 \(S\)-群同态。下列条件等价：}

\begin{enumerate}\item[(i)]\SGARefTarget{sga3:xi:corollary:5:4:item:i}\textit{\(u_1\) 和 \(u_2\) 对平展拓扑局部共轭。}

\item[(i bis)]\SGARefTarget{sga3:xi:corollary:5:4:item:i-bis}\textit{\(u_1\) 和 \(u_2\) 对忠实平坦拟紧拓扑局部共轭。}

\item[(ii)]\SGARefTarget{sga3:xi:corollary:5:4:item:ii}\textit{对每个 \(s\in S\)，以 \(\overline{s}\) 表示 \(\kappa(s)\) 的代数闭包的谱。态射 \[
  u_{1,\overline{s}},u_{2,\overline{s}}:
  H_{\overline{s}}\rightrightarrows G_{\overline{s}}
\] 在 \(G\bigl(\kappa(\overline{s})\bigr)\) 的某个元素作用下共轭。}

\item[(ii bis)]\SGARefTarget{sga3:xi:corollary:5:4:item:ii-bis}\textit{结构态射 \(\operatorname{Transp}(u_1,u_2)\to S\) 为满射。}

\item[(iii)]\SGARefTarget{sga3:xi:corollary:5:4:item:iii}\textit{\(\operatorname{Transp}(u_1,u_2)\) 为有限型平滑 \(S\)-群预概形 \(\operatorname{Centr}_G(u_1)\) 作用下的挠子。}\end{enumerate}

蕴含关系 \SGARefLink{sga3:xi:corollary:5:4:item:i}{(i)} \(\Rightarrow\) \SGARefLink{sga3:xi:corollary:5:4:item:i-bis}{(i bis)} 和 \SGARefLink{sga3:xi:corollary:5:4:item:ii}{(ii)} \(\Rightarrow\) \SGARefLink{sga3:xi:corollary:5:4:item:ii-bis}{(ii bis)} 直接成立，后者由零点定理得到。蕴含关系 \SGARefLink{sga3:xi:corollary:5:4:item:i-bis}{(i bis)} \(\Rightarrow\) \SGARefLink{sga3:xi:corollary:5:4:item:ii}{(ii)} 由“有限扩张原理”（EGA IV，9）得出。此外，\SGARefLink{sga3:xi:corollary:5:4:item:ii-bis}{(ii bis)} \(\Rightarrow\) \SGARefLink{sga3:xi:corollary:5:4:item:iii}{(iii)}：由于 \(\operatorname{Transp}(u_1,u_2)\) 在 \(S\) 上平滑，它在 \(S\) 上平坦；又由于它有限型，故在 \(S\) 上拟紧。因此它在 \(S\) 上忠实平坦且拟紧，当且仅当其结构态射满射。另一方面，它在 \(\operatorname{Centr}_G(u_1)\) 作用下形式主齐性，而后者在 \(S\) 上忠实平坦且拟紧。因此最后一个条件也等价于说，在忠实平坦拟紧拓扑中，\(\operatorname{Transp}(u_1,u_2)\) 为 \(\operatorname{Centr}_G(u_1)\) 作用下的挠子。最后，\SGARefLink{sga3:xi:corollary:5:4:item:iii}{(iii)} \(\Rightarrow\) \SGARefLink{sga3:xi:corollary:5:4:item:i}{(i)} 由 \(\operatorname{Transp}(u_1,u_2)\) 在 \(S\) 上平滑以及形式为 \SGARefLink{sga3:xi:proposition:1:10}{1.10} 的 Hensel 引理推出。

\medskip\noindent\SGARefTarget{sga3:xi:remark:5:5}\textbf{注 5.5。} --- 固定 \(u_1=u:H\to G\)。考虑函子 \SGARefTarget{sga3:xi:display:5:5:functor} \[
  (\mathbf{Sch})^\circ/S\longrightarrow(\mathbf{Ens})
\]，它将 \(S\) 上的每个 \(T\) 映为所有对平展拓扑局部共轭于 \(u_T:H_T\to G_T\) 的 \(T\)-群同态 \(u_2:H_T\to G_T\) 的集合。\SGARefLink{sga3:xi:display:5:5:functor}{此函子} 恰由 \SGARefLink{sga3:xi:corollary:5:3}{5.3} 中考虑的、同构于 \(G/\operatorname{Centr}_G(u)\) 的 \(M\) 的开子预概形表示。

简述以 \SGARefLink{sga3:xi:theorem:4:1}{4.1} 代替 \SGARefLink{sga3:xi:corollary:4:2}{4.2} 所得到的上述结果的变体。设 \(G\) 为 \(S\) 上的平滑仿射群预概形，现在以 \(M\) 表示 \(S\) 上分离的平滑 \(S\)-预概形，它表示 \SGARefLink{sga3:xi:theorem:4:1}{4.1} 中考虑的函子。\(G\) 仍作用于 \(M\)：\SGARefTarget{sga3:xi:display:5:subgroup-conjugation-action} \[
  (g,H)\longmapsto\operatorname{int}(g)(H).
\]。因此，与上面一样，存在态射 \SGARefTarget{sga3:xi:display:5:x-bis} \[
  \tag{x bis}
  G\times_S M\longrightarrow M\times_S M.
\]。

\medskip\noindent\SGARefTarget{sga3:xi:corollary:5:1:bis}\textbf{推论 5.1 bis。} --- \textit{\SGARefLink{sga3:xi:display:5:x-bis}{上述态射} 平滑。}

这由 \SGARefLink{sga3:xi:theorem:4:1}{4.1} 和 \SGARefLink{sga3:xi:corollary:2:3:bis}{2.3 bis} 得出。再次得到：

\medskip\noindent\SGARefTarget{sga3:xi:corollary:5:2:bis}\textbf{推论 5.2 bis。} --- \textit{设 \(H_1,H_2\) 为 \(G\) 的两个乘法型子群。则 \(G\) 的子函子 \(\operatorname{Transp}_G(H_1,H_2)\) 可由 \(G\) 的闭子预概形表示，且在 \(S\) 上平滑。}

特别地：

\medskip\noindent\SGARefTarget{sga3:xi:corollary:5:3:bis}\textbf{推论 5.3 bis。} --- \textit{设 \(H\) 为 \(G\) 的乘法型子群。则 \(G\) 的群子函子 \(\operatorname{Norm}_G(H)\) 可由 \(G\) 的闭子预概形表示，且在 \(S\) 上平滑。此外，商 \(G/\operatorname{Norm}_G(H)\) 可由 \(M\) 的开子预概形表示。}

\medskip\noindent\SGARefTarget{sga3:xi:corollary:5:4:bis}\textbf{推论 5.4 bis。} --- \textit{设 \(H_1,H_2\) 为 \(G\) 的两个乘法型子群。下列条件等价：}

\begin{enumerate}\item[(i)]\SGARefTarget{sga3:xi:corollary:5:4:bis:item:i}\textit{\(H_1\) 和 \(H_2\) 对平展拓扑局部共轭。}

\item[(i bis)]\SGARefTarget{sga3:xi:corollary:5:4:bis:item:i-bis}\textit{\(H_1\) 和 \(H_2\) 对忠实平坦拟紧拓扑局部共轭。}

\item[(ii)]\SGARefTarget{sga3:xi:corollary:5:4:bis:item:ii}\textit{对每个 \(s\in S\)，以 \(\overline{s}\) 表示 \(\kappa(s)\) 的代数闭包的谱。\(G_{\overline{s}}\) 的子群 \(H_{1,\overline{s}}\) 和 \(H_{2,\overline{s}}\) 在 \(G\bigl(\kappa(\overline{s})\bigr)\) 的某个元素作用下共轭。}

\item[(ii bis)]\SGARefTarget{sga3:xi:corollary:5:4:bis:item:ii-bis}\textit{结构态射 \(\operatorname{Transp}(H_1,H_2)\to S\) 为满射。}

\item[(iii)]\SGARefTarget{sga3:xi:corollary:5:4:bis:item:iii}\textit{\(\operatorname{Transp}(H_1,H_2)\) 为有限型平滑 \(S\)-群预概形 \(\operatorname{Norm}_G(H_1)\) 作用下的主齐性丛。}\end{enumerate}

\medskip\noindent\SGARefTarget{sga3:xi:remark:5:5:bis}\textbf{注 5.5 bis。} --- \SGARefLink{sga3:xi:remark:5:5}{注 5.5} 同样适用于当前情形。

\medskip\noindent\SGARefTarget{sga3:xi:remark:5:6}\textbf{注 5.6。} --- 为证明 \SGARefLink{sga3:xi:corollary:5:2}{5.2}，从而证明 \SGARefLink{sga3:xi:corollary:5:3}{5.3} 的第一个断言以及 \SGARefLink{sga3:xi:corollary:5:4}{5.4} 的各断言，可以将对 \SGARefLink{sga3:xi:corollary:4:2}{4.2} 的引用换为对 VIII，\SGARefLink{sga3:viii:theorem:6:4}{6.4} 的引用，后者的证明简单得多。这还表明，在那里无需假设 \(G\) 在 \(S\) 上仿射。此外，\SGARefLink{sga3:xi:section:6}{第 6 节} 将说明这种方法的一种变体如何将这些结果推广到某些比乘法型群更一般的群 \(H\)。同样的观察适用于变体 \SGARefLink{sga3:xi:corollary:5:2:bis}{5.2 bis} 及其后各条。反之，下面的结果 \SGARefLink{sga3:xi:theorem:5:8}{5.8} 本质地使用了所有假设，特别是 \(G\) 在 \(S\) 上仿射且平滑、\(H\) 为乘法型群；报告者不知道除使用可表性定理 \SGARefLink{sga3:xi:theorem:4:1}{4.1} 或 \SGARefLink{sga3:xi:corollary:4:2}{4.2} 外的证明。\footnote{自本文写成以来，情况已有变化；参见第 XV、XIX 讲第~6 节。}

\medskip\noindent\SGARefTarget{sga3:xi:paragraph:5:7}\textbf{5.7。}由于态射 \(\mathrm{\SGARefLink{sga3:xi:display:5:x}{(x)}}\)（分别为 \SGARefLink{sga3:xi:display:5:x-bis}{\(\mathrm{(x\ bis)}\)}）平滑，因而为开态射，其像是开的。设 \(R\) 为此像，赋予从 \(M\times_S M\) 诱导的结构。容易验证，\(R\) 为 \(M\) 上的等价关系，即 \SGARefLink{sga3:xi:corollary:5:4}{5.4}（分别为 \SGARefLink{sga3:xi:corollary:5:4:bis}{5.4 bis}）中明确给出的关系。值得研究的是，在 \(S\) 上形式平展的商层 \(M/R\) 是否可表；若可表，则它由 \(S\) 上平展的预概形表示。例如，当 \(S\) 为域的谱时就是如此。还需注意，\(R\) 未必在 \(M\times_S M\) 中闭，即便 \(G\) 在 \(S\) 上拟有限也如此。因此，当 \(M/R\) 可表时，它在 \(S\) 上可能不分离。

\medskip\noindent\SGARefTarget{sga3:xi:theorem:5:8}\textbf{定理 5.8。} --- \textit{设 \(S\) 为预概形，\(G\) 为平滑仿射 \(S\)-预概形，且 \(s\in S\)。则：}

\begin{enumerate}\item[(a)]\SGARefTarget{sga3:xi:theorem:5:8:item:a}\textit{对 \(G_s\) 的每个乘法型子群 \(H_0\)，存在平展态射 \(S'\to S\)、\(s\) 上方的点 \(s'\in S'\)，使剩余域扩张 \(\kappa(s')/\kappa(s)\) 平凡，以及 \(G'=G\times_S S'\) 的乘法型子群 \(H'\)，使得} \SGARefTarget{sga3:xi:display:5:8:item:a:fiber} \[
  H'_{s'}=H_0\otimes_{\kappa(s)}\kappa(s').
\]。

\item[(b)]\SGARefTarget{sga3:xi:theorem:5:8:item:b}\textit{对每个群同态 \(u_0:H_s\to G_s\)，其中 \(H\) 为有限型乘法型 \(S\)-群预概形，存在平展态射 \(S'\to S\)、\(s\) 上方的点 \(s'\in S'\)，使 \(\kappa(s')/\kappa(s)\) 平凡，以及群同态 \[
  u':H'\longrightarrow G',
  \qquad
  H'=H\times_S S',
  \quad
  G'=G\times_S S',
\]，使得} \SGARefTarget{sga3:xi:display:5:8:item:b:fiber} \[
  u'_{s'}:H'_{s'}\longrightarrow G'_{s'}
  \quad\text{等于}\quad
  u_0\otimes_{\kappa(s)}\kappa(s').
\]。\end{enumerate}

这由 \SGARefLink{sga3:xi:theorem:4:1}{4.1}（分别为 \SGARefLink{sga3:xi:corollary:4:2}{4.2}）及形式为 \SGARefLink{sga3:xi:proposition:1:10}{1.10} 的 Hensel 引理得到。

\medskip\noindent\SGARefTarget{sga3:xi:proposition:5:9}\textbf{命题 5.9。} --- \textit{设 \(S\) 为预概形，\(G\) 为平滑仿射 \(S\)-预概形，\(H\) 为 \(G\) 的乘法型子群。考虑} \[
  \operatorname{Centr}_G(H)\lhook\joinrel\longrightarrow
  \operatorname{Norm}_G(H),
\]\textit{，由 \SGARefLink{sga3:xi:corollary:5:3}{5.3} 和 \SGARefLink{sga3:xi:corollary:5:3:bis}{5.3 bis}，它们是 \(G\) 的闭群子预概形，在 \(S\) 上平滑。则前者是后者的开闭子预概形，商层} \SGARefTarget{sga3:xi:display:5:9:weyl-quotient} \[
  W_G(H)=
  \operatorname{Norm}_G(H)\big/\operatorname{Centr}_G(H)
\]\textit{由 \(\operatorname{Aut}_{S\text{-}\mathrm{gr}}(H)\) 的开群子预概形表示。因此它是拟有限、平展且分离的 \(S\)-群预概形。}

事实上，考虑显然的同态 \SGARefTarget{sga3:xi:display:5:9:theta} \[
  \theta:\operatorname{Norm}_G(H)\longrightarrow
  \operatorname{Aut}_{S\text{-}\mathrm{gr}}(H),
\]，其核按定义为 \(\operatorname{Centr}_G(H)\)。由于 \(\operatorname{Aut}_{S\text{-}\mathrm{gr}}(H)\) 由 \(S\) 上分离的平展 \(S\)-群预概形表示（X，\SGARefLink{sga3:x:corollary:5-10}{5.10}），其单位截面为开闭浸入。因此它在 \(\theta\) 下的逆像为 \(\operatorname{Norm}_G(H)\) 的开闭子群。

此外，\(\theta\) 平滑。这形式地由定义以及 \(\operatorname{Norm}_G(H)\) 在 \(S\) 上平滑、\(\operatorname{Aut}_{S\text{-}\mathrm{gr}}(H)\) 在 \(S\) 上平展推出。与 \SGARefLink{sga3:xi:corollary:5:3}{5.3} 中一样，\(\theta\) 的像因而为 \(\operatorname{Aut}_{S\text{-}\mathrm{gr}}(H)\) 的开子预概形 \(U\)，且赋予诱导结构后，\(U\) 表示商层 \(\operatorname{Norm}_G(H)/\operatorname{Centr}_G(H)\)。所以后者在 \(S\) 上平展且分离。它还在 \(S\) 上拟有限，因为它在 \(S\) 上拟紧，这是由于它是拟紧的 \(\operatorname{Norm}_G(H)\) 的像。这就证明了 \SGARefLink{sga3:xi:proposition:5:9}{5.9}。

\medskip\noindent\SGARefTarget{sga3:xi:corollary:5:10}\textbf{推论 5.10。} --- \textit{对每个 \(s\in S\)，置} \SGARefTarget{sga3:xi:display:5:10:w} \[
  w(s)=
  \operatorname{rank}
  \bigl(
    \operatorname{Norm}_{G_s}(H_s)/
    \operatorname{Centr}_{G_s}(H_s)
  \bigr).
\]\textit{。这也是 \(\operatorname{Centr}_{G(\overline{k})}(H(\overline{k}))\) 在 \(\operatorname{Norm}_{G(\overline{k})}(H(\overline{k}))\) 中的指数，其中 \(\overline{k}\) 为 \(\kappa(s)\) 的代数闭包。则函数 \(s\mapsto w(s)\) 下半连续。它在 \(s\) 的某个邻域内为常值，当且仅当 \(W_G(H)\) 在 \(s\) 的某个邻域内于 \(S\) 上有限。}

事实上，对每个在 \(S\) 上平展、有限型且分离的 \(S\)-预概形 \(W\)，函数 \[
  s\longmapsto
  \operatorname{rank}[W_s:\kappa(s)]
\] 下半连续，且在 \(s\) 的某个邻域内为常值，当且仅当 \(W\) 在 \(s\) 的某个邻域内于 \(S\) 上有限。这一事实曾在 SGA 1，I，10.9 中指出；其证明并不困难，见 EGA IV。\footnote{EGA IV，15.5.1 和 18.10.7。}

\medskip\noindent\SGARefTarget{sga3:xi:remark:5:11}\textbf{注 5.11。}\footnote{编者注：向非仿射情形的推广见 M.~Raynaud，\textit{群概形与齐性空间上的丰沛层〔Faisceaux amples sur les schémas en groupes et les espaces homogènes〕}，Lecture Notes in Mathematics 119 (1970)，IX.2.8。} --- 设 \(G\) 为 \(S\) 上仿射且平滑的群预概形，\(H\) 为乘法型子群。则 \SGARefLink{sga3:xi:corollary:5:3}{5.3} 和 \SGARefLink{sga3:xi:corollary:5:3:bis}{5.3 bis} 蕴含商 \SGARefTarget{sga3:xi:display:5:11:quotients} \[
  G/\operatorname{Centr}_G(H)
  \qquad\text{且}\qquad
  G/\operatorname{Norm}_G(H)
\] 是 \(S\) 上平滑且拟仿射的预概形。事实上，在两种情形中，商所嵌入的模预概形 \(M\) 都具有下述性质：每个在 \(S\) 上拟紧的开集 \(U\subset M\)，在 \(S\) 上拟仿射；这是由 \SGARefLink{sga3:xi:corollary:3:11}{3.11} 得到的。

此外，我们将在下一讲看到\footnote{编者注：参见注 XII.5.7。这里提及如下推广：不假设 \(G\) 仿射，若 \(H\) 为环面，\(S\) 局部 Noether，则由 Raynaud 定理，\(U\) 仿射且平滑（\textit{同处}，IX.2.6 和 IX.2.8）。事实上，\(G/\operatorname{Norm}_G(H)\) 为 \(M\) 的开子预概形，而 \(M\) 的最大可表开子函子是在 \(S\) 上平滑且仿射的开子预概形的不交和。}，对每个这样的 \(U\)，\(U\) 在 \(M\) 中的概形闭包 \(\overline{U}\) 甚至在 \(S\) 上仿射。因此，只要相应模预概形（即从 \(H\) 到 \(G\) 的群同态、分别为 \(G\) 的乘法型子群的模预概形）的开子预概形 \(U\) 在 \(M\) 中闭，上述商就在 \(S\) 上仿射。例如，第二种情形中 \(H\) 为“极大环面”时，或 \(S\) 为域的谱时，就是如此；参见 XII.5.4。\footnote{编者注：此断言假设 \(G\) 具有局部常值的约化秩。}我不知道上述商是否一般都在 \(S\) 上仿射。


